\section{Radiation scheme} \label{sec:basic_usel_radiation}
%-------------------------------------------------------------------------------
The radiation process is configured by two namelists.
First, the radiation component is enabled by \nmitem{ATMOS_PHY_RD_DO} in \namelist{PARAM_ATMOS}.
Then, the radiation scheme and its calling interval are specified in \namelist{PARAM_ATMOS_PHY_RD}.

\editboxtwo{
 \verb|&PARAM_ATMOS | & \\
 \verb|ATMOS_PHY_RD_DO = .true.,| & ; Enable radiation \\
 \verb|/ | & \\
}
\editboxtwo{
 \verb|&PARAM_ATMOS_PHY_RD | & \\
 \verb|RD_TYPE          = "SIMPLE",| & ; Choose from radiation schemes in Table \ref{tab:rd_scheme_list} \\
 \verb|RD_SOLARINS_TYPE = "SIMPLE",| & ; Type of incoming solar-insolation scheme \\ 
 \verb|TIME_DT      = 600.D0,|      & ; Time interval of the radiation process \\
 \verb|TIME_DT_UNIT = "SEC",|      & ; Unit for \verb|TIME_DT| \\
 \verb|/ | & \\
}

The available choices of \nmitem{RD_TYPE} are listed in Table \ref{tab:rd_scheme_list}.
The incoming solar-insolation scheme is selected by \nmitem{RD_SOLARINS_TYPE}.
Currently, \verb|"SIMPLE"| and \verb|"NONE"| are available.
When \verb|"SIMPLE"| is selected, its detailed configuration is given in \namelist{PARAM_ATMOS_SOLARINS_SIMPLE} (Section~\ref{subsec:physics_rd_solar_radiation}).

%%%%%%%%%%%
\begin{table}[htbp]
\caption{Radiation schemes available in \scaledg.}
\label{tab:rd_scheme_list}
\begin{center}
\begin{tabularx}{150mm}{lXX}
\hline
\rowcolor[gray]{0.9}
\verb|RD_TYPE| & Description & Reference \\
\hline
\verb|SIMPLE|  & Simplified radiation scheme & \cite{Geen2016,Vallis2018} \\
\hline
\end{tabularx}
\end{center}
\end{table}
%%%%%%%%%%%

%------------------------------------------------------------------------
\subsection{Configuration of incoming solar radiation}
\label{subsec:physics_rd_solar_radiation}

The incoming solar radiation used by the simplified radiation scheme is prescribed by an idealized solar-insolation model.
Two types of idealized forcing are currently available: a spatially uniform constant flux and an annual-mean insolation that varies with latitude.
These settings are configured in \namelist{PARAM_ATMOS_SOLARINS_SIMPLE}.


The choice of \nmitem{SOLARINS_TYPE} determines how the incoming solar flux is specified.
%%
\editboxtwo{
\verb|&PARAM_ATMOS_SOLARINS_SIMPLE | & \\
\verb|SOLARINS_TYPE          = "ANNUAL_MEAN",| & ; Type of idealized solar insolation; \\
                                               & ;   ANNUAL\_MEAN or CONST \\
\verb|ORBIT_REFERENCE_YEAR   = 2000,|           & ; Reference year for orbital parameters \\
\verb|ANNUAL_MEAN_YEAR       = 2000,|           & ; Year used for the annual-mean calculation \\
\verb|ANNUAL_SAMPLES_PER_DAY = 1,|              & ; Number of samples per day \\
\verb|/ | & \\
}
When \verb|SOLARINS_TYPE="ANNUAL_MEAN"| is selected, the daily-mean solar insolation is calculated from the orbital configuration and latitude and then averaged over \nmitem{ANNUAL_MEAN_YEAR}.
The default is \verb|"ANNUAL_MEAN"|.
\nmitem{ORBIT_REFERENCE_YEAR} specifies the reference year used to initialize the orbital parameters in the SCALE solar-insolation module.
\nmitem{ANNUAL_SAMPLES_PER_DAY} specifies the number of equally spaced samples used within each day when forming the annual mean.
The default value is one sample per day.


When \verb|SOLARINS_TYPE="CONST"| is selected, \nmitem{CONST_FLUX} is applied uniformly over the model domain.
%%
\editboxtwo{
\verb|&PARAM_ATMOS_SOLARINS_SIMPLE | & \\
\verb|SOLARINS_TYPE          = "CONST",| & ; Type of idealized solar insolation; \\
\verb|CONST_FLUX             = 340.D0,|         & ; Constant incoming solar flux [W m$^{-2}$] \\
\verb|/ | & \\
}


%------------------------------------------------------------------------
\subsection{Simplified radiation scheme}
\label{subsec:physics_rd_simple}
%%

The simplified radiation schemes are intended for idealized atmospheric experiments.
Compared with comprehensive radiative transfer models, they employ simplified spectral representations 
while retaining the dominant radiative heating processes.

Two formulations are available: a gray longwave radiation scheme and a two-band longwave plus one-band shortwave radiation scheme.
The formulation and its optical-depth parameters are configured in \namelist{PARAM_ATMOS_PHY_RD_DGM_SIMPLE}.  
The default is the gray longwave radiation scheme.


\subsubsection{Gray longwave radiation scheme}
A basic configuration using the default parameters is, for example,
%%
\editboxtwo{
\verb|&PARAM_ATMOS_PHY_RD_DGM_SIMPLE | & \\
\verb|SIMPLE_RD_TYPE                  = "Gray",|  & ; Type of simplified radiation scheme \\
\verb|V2018_GRAY_OPTDEPTH_PARAMS_TYPE = "V2018",| & ; Parameter set for gray optical depth \\
\verb|DiffusivityFactor               = 1.D0,|    & ; Diffusivity factor for two-stream approximation \\
\verb|pCO2                            = 360.D0,|  & ; CO$_2$ concentration [ppmv] \\
\verb|/ | & \\
}

For the gray scheme, the optical-depth parameters are controlled by \nmitem{V2018_GRAY_OPTDEPTH_PARAMS_TYPE}.
The available choices are \verb|"BO2013"|, \verb|"V2018"|, and \verb|"USER"|, with \verb|"V2018"| as the default.
When \verb|"USER"| is selected, the four coefficients \nmitem{V2018_GRAY_OPTDEPTH_LW_PARAMS} must be specified in the order $A,\mu,B,C$, 
corresponding to Eq.~(5) of \cite{Vallis2018}.

The preset coefficient values used internally are $(A,\mu,B,C)=(0.8678,1.0,1997.9,0.0)$ for \verb|"BO2013"| and
$(0.1627,1.0,1997.9,0.17)$ for \verb|"V2018"|.

For the gray radiation scheme, the shortwave optical depth is set to zero.
Therefore, the prescribed incoming solar flux reaches the surface without attenuation by atmospheric shortwave absorption.


\subsubsection{Two-band longwave and one-band shortwave scheme}

\editboxtwo{
\verb|&PARAM_ATMOS_PHY_RD_DGM_SIMPLE | & \\
\verb|SIMPLE_RD_TYPE = "LWbnd2+SWbnd1",| & ; Two LW bands and one SW band \\
\verb|V2018_LWbnd2SWbnd1_OPTDEPTH_PARAMS_TYPE = "V2018",| & ; Optical-depth parameter set \\
\verb|V2018_LWbnd2SWbnd1_OPTDEPTH_SW_kadd = 4,| & ; Number of layers added above the model top \\
\verb|V2018_LWbnd2SWbnd1_OPTDEPTH_SW_QV_above_MOT = 1.D-3,| & ; Specific humidity above the model top \\
\verb|DiffusivityFactor = 1.D0,| & ; Diffusivity factor \\
\verb|pCO2              = 360.D0,| & ; CO$_2$ concentration [ppmv] \\
\verb|/ | & \\
}

For this scheme, \nmitem{V2018_LWbnd2SWbnd1_OPTDEPTH_PARAMS_TYPE} can be set to \verb|"V2018"| or \verb|"USER"|.
When \verb|"V2018"| is selected, the optical-depth coefficients and the longwave spectral fractions are prescribed internally following \cite{Vallis2018}.
When \verb|"USER"| is selected, the user specifies
\nmitem{V2018_LWbnd2SWbnd1_OPTDEPTH_{SW,LW,LW_window}_PARAMS} 
and \nmitem{V2018_LWbnd2SWbnd1_LW_window_FRAC}. 

The shortwave coefficients contain $a_{\rm SW}$ and $c_{\rm SW}$;
the non-window longwave coefficients contain $a_{\rm LW}$, $b_{\rm LW}$, $c_{\rm LW}$, $d_{\rm LW}$;
and the longwave-window coefficients contain $a_{\rm win}$, $b_{\rm win}$, $c_{\rm win}$, $d_{\rm win}$.
The window fraction specifies the fraction of the longwave spectrum assigned to the $8-14~\mu\mathrm{m}$ window band.

The vertical coordinate in SCALE-DG is based on the height-based terrain-following coordinate.
Thus, the optical depth above the model top is not well defined.
To address this, the two parameters
\nmitem{V2018_LWbnd2SWbnd1_OPTDEPTH_SW_kadd} and
\nmitem{V2018_LWbnd2SWbnd1_OPTDEPTH_SW_QV_above_MOT} control the treatment of shortwave optical thickness above the model top; 
their default values are $4$ and $10^{-3}$, respectively.

\subsubsection{Common parameters for simplified radiation schemes}
\nmitem{DiffusivityFactor} is the diffusivity factor used in the two-stream approximation. Although a value around 1.66 is commonly used for a plane-parallel atmosphere with isotropic scattering, 
the present idealized scheme uses 1.0 following \cite{Vallis2018}.
\nmitem{pCO2} specifies the CO$_2$ concentration in ppmv and has a default value of 360.  
